

\section*{Chapitre \ref{ch:iv3:python-and-pandas} --- Python et pandas}
\markboth{\protect\omTocMark{}{Réponses~: Python et pandas}}{\protect\omTocMark{}{Réponses~: Python et pandas}}

\begin{solution}{iq:iv3:python-and-pandas:1}
Il affiche \texttt{[1]}, \texttt{[1, 2]}, \texttt{[3]} et \texttt{[1, 2, 4]}. Une valeur par défaut est évaluée une seule fois, au moment où l'instruction \texttt{def} s'exécute, si bien que chaque appel qui omet la liste partage la même liste, laquelle survit à l'appel.

Le troisième appel passe sa propre liste~; il ne touche donc pas à la liste partagée, ce qui explique que le quatrième poursuive à partir de \texttt{[1, 2]}.

La correction consiste en la sentinelle habituelle~: on écrit \texttt{items=None} dans la signature et, à l'intérieur de la fonction, \texttt{if items is None: items = []}. Une nouvelle liste est alors créée à chaque appel. On teste \texttt{is None} et non \texttt{if not items}, car un appelant qui passe une liste vide attend qu'on lui renvoie cette même liste.

\iqfollow{Où la liste est-elle conservée~?}{Sur l'objet fonction~: \texttt{add.\_\_defaults\_\_} la contient, et on peut la voir grandir.}
\iqfollow{Cela peut-il avoir un intérêt~?}{Parfois comme cache peu coûteux, mais un lecteur y verra un bogue~; on utilise donc plutôt \texttt{functools.cache}.}
\end{solution}

\begin{solution}{iq:iv3:python-and-pandas:2}
Il affiche \texttt{\{\textquotesingle A\textquotesingle: [1], \textquotesingle B\textquotesingle: [1]\}} et \texttt{[[9, 0], [9, 0]]}. Dans les deux cas, un seul objet liste est partagé, si bien qu'une modification faite sous un nom se voit sous tous les autres.

\texttt{dict.fromkeys} évalue son second argument une seule fois et range cet unique objet sous chaque clé. \texttt{[[0] * 2] * 2} répète une référence à la même ligne interne. Python ne copie pas lors de la répétition, et il ne copie jamais lors d'une affectation.

Les corrections créent une liste par clé ou par ligne~: \texttt{\{k: [] for k in names\}}, ou un \texttt{defaultdict(list)}, et \texttt{[[0] * 2 for \_ in range(2)]}, où l'expression est évaluée une fois par ligne.

Le 0 est sans danger, car un entier ne peut pas changer. \texttt{d[k] += 1} rattache cette seule clé à un nouvel entier~; il ne modifie pas le 0 partagé. Le partage est donc sans effet néfaste pour les valeurs immuables, et l'expression interne \texttt{[0] * 2} de la ligne ci-dessus est correcte pour la même raison.

\iqwrong{«~\texttt{[0] * 2} est le bogue.~» Ce n'est pas le cas, puisque les éléments sont des entiers. C'est le \texttt{* 2} externe qui répète une liste.}
\iqfollow{Comment vérifiez-vous que deux noms désignent un même objet~?}{\texttt{a is b}, ou comparer \texttt{id(a)} et \texttt{id(b)}.}
\end{solution}

\begin{solution}{iq:iv3:python-and-pandas:3}
Il affiche \texttt{True True}, puis \texttt{False True}. \texttt{==} compare les valeurs, et \texttt{is} compare les identités~: il indique si les deux noms désignent exactement le même objet.

CPython conserve un objet partagé pour chaque entier de $-5$ à 256, si bien que \texttt{int(\dots)} renvoie l'objet mis en cache pour 256, et un nouvel objet à chaque fois pour 257. Il s'agit d'un détail d'implémentation, que d'autres implémentations ou versions n'ont pas à suivre.

Le code passe volontairement par \texttt{int} sur un texte. Avec les littéraux \texttt{257 is 257} dans un même bloc, le compilateur peut fusionner les constantes, et l'on obtient \texttt{True}, ce qui ajoute à la confusion.

On utilise \texttt{is} pour les singletons, où c'est l'identité qui compte~: \texttt{x is None}, ou un objet sentinelle propre. On utilise \texttt{==} pour tout le reste, y compris \texttt{True} et \texttt{False} venus d'une bibliothèque, où un booléen NumPy n'est pas l'objet \texttt{True}.

\iqwrong{«~\texttt{is} est un \texttt{==} plus rapide.~» Il est plus rapide, et il répond à une autre question. Du code qui compare des quantités avec \texttt{is} passe tous les petits tests, puis échoue dès 257 actions.}
\end{solution}

\begin{solution}{iq:iv3:python-and-pandas:4}
Il affiche \texttt{tuple}, puis une \texttt{TypeError} pour la liste, une \texttt{TypeError} pour le tuple qui contient une liste, et \texttt{absent} pour le frozenset. Une clé doit être hachable, et une liste ne l'est pas.

Un dictionnaire range une clé selon sa valeur de hachage, et il retrouve la clé en recalculant cette valeur. Si une clé pouvait changer après insertion, sa valeur de hachage changerait, et l'entrée se trouverait au mauvais endroit, inaccessible. C'est pourquoi les types natifs modifiables refusent d'être hachés. La vérification intervient à toute recherche, puisque \texttt{get} doit aussi hacher son argument.

La valeur de hachage d'un tuple se construit à partir de ses éléments~; un tuple n'est donc hachable que si tous ses éléments le sont. \texttt{(1, [2])} échoue pour cette raison, et la vérification ne porte pas sur le fait d'être «~un tuple~».

Pour la mémoire, un tuple est plus petit, car il n'a pas besoin de place libre pour grandir. Sur CPython 3.10 en 64 bits, \texttt{sys.getsizeof} donne 56 pour un tuple de deux éléments et 72 pour une liste de deux éléments écrite entre crochets. La différence est faible pour un seul objet, et importante pour un million de petits enregistrements.

\iqfollow{Le contenu d'un tuple peut-il changer~?}{On ne peut pas réaffecter le tuple à d'autres objets, mais une liste qu'il contient peut toujours être modifiée. Donc «~immuable~» n'est vrai qu'en surface.}
\end{solution}

\begin{solution}{iq:iv3:python-and-pandas:5}
Il affiche \texttt{\{\textquotesingle A\textquotesingle: [1, 2, 3]\}} et \texttt{\{\textquotesingle A\textquotesingle: [1, 2]\}}. La copie superficielle est un nouveau dictionnaire qui contient les mêmes listes~; elle ne voit donc pas la nouvelle clé \texttt{B}, mais elle voit l'ajout. La copie profonde possède de nouvelles listes à tous les niveaux.

Pour l'instantané, \texttt{deepcopy} parcourt toute la structure avec une table de mémoïsation, à un coût au niveau de Python pour chaque objet~; appelée plusieurs fois par seconde sur un gros carnet d'ordres, elle apparaît dans un profilage.

Les options préférables dépendent de ce qui est lu. Si un instantané n'est destiné qu'à un seul consommateur, on envoie une forme compacte~: un tuple de couples \texttt{(price, size)} pour chaque côté du carnet. Les tuples de valeurs immuables peuvent être partagés sans risque~; la copie du seul premier niveau suffit alors. Ou bien on conserve les niveaux comme valeurs immuables, on en remplace un à chaque modification, et l'on prend la copie superficielle en $O(\text{niveaux})$.

\iqwrong{«~\texttt{dict(orders)} la rend indépendante.~» C'est la même copie superficielle. Cela paraît correct dans un test où l'on ne fait qu'ajouter des clés.}
\iqfollow{Et pour les tableaux NumPy~?}{\texttt{arr.copy()} est une copie profonde pour les types numériques, et elle copie un seul tampon en un seul appel, sans étape au niveau de Python pour chaque élément.}
\end{solution}

\begin{solution}{iq:iv3:python-and-pandas:6}
Il affiche \texttt{False True False}, puis \texttt{nan 1.0}, puis \texttt{[2.0, nan, 1.0]}. Un NaN donne faux à toute comparaison, lui-même compris, et cette seule règle explique les trois lignes.

\texttt{nan == nan} est faux. \texttt{nan in [nan]} est vrai, car \texttt{in} teste d'abord l'identité, et l'objet est le même, alors qu'un \texttt{float(\textquotesingle nan\textquotesingle)} neuf est un autre objet, et l'égalité échoue.

\texttt{max} garde son premier argument, sauf si un argument suivant est strictement supérieur. Rien n'est strictement supérieur à NaN, donc \texttt{max(nan, 1.0)} reste à NaN~; et NaN n'est pas supérieur à 1.0, donc \texttt{max(1.0, nan)} reste à 1.0. Le résultat dépend donc de la position du NaN.

\texttt{sorted} traite toute comparaison avec NaN comme fausse~; la liste est donc considérée comme déjà triée, et elle reste non triée, sans erreur. La liste donne alors un minimum, un maximum et une médiane erronés.

Pour le prix le plus élevé, on filtre d'abord~: \texttt{max(x for x in prices if not math.isnan(x))}, ou \texttt{np.nanmax}. Il ne faut pas appeler \texttt{max} en comptant sur la chance.

\iqwrong{«~Tester \texttt{x == float(\textquotesingle nan\textquotesingle)}.~» Ce test n'est jamais vrai. On utilise \texttt{math.isnan}.}
\end{solution}

\begin{solution}{iq:iv3:python-and-pandas:7}
Il affiche \texttt{[(1, 3), (2, 4)]}, \texttt{[]}, \texttt{5 0} et \texttt{3 3}. \texttt{zip} et un générateur sont des itérateurs, que l'on ne peut parcourir qu'une seule fois. Un \texttt{range} est une séquence, qui ouvre un nouvel itérateur à chaque fois.

Ainsi, sur un générateur de 1, 2 et 3, la fonction compte 3, ce qui épuise le flux, puis ne somme rien et renvoie $0 / 3 = 0.0$, une moyenne plausible. Si la somme est calculée en premier, elle obtient la bonne somme et une longueur de 0, et lève une \texttt{ZeroDivisionError}, ce qui est l'échec le moins trompeur.

La correction dépend de la taille. Si les données tiennent en mémoire, \texttt{xs = list(xs)} en tête de fonction. Sinon, il faut une seule passe qui tient à jour un compteur et une somme courante. Le test \texttt{iter(xs) is xs} indique si l'on a affaire à un itérateur à usage unique, ce qu'une fonction peut vouloir refuser.

\iqwrong{«~Utiliser \texttt{itertools.tee}.~» Cela fonctionne, mais il stocke tout ce qu'une branche a lu et que l'autre n'a pas lu, soit ici tout le flux, si bien qu'il n'économise rien.}
\iqfollow{Pourquoi \texttt{for} ne lève-t-il rien à la fin~?}{Un itérateur épuisé lève aussitôt une \texttt{StopIteration} à chaque demande ultérieure, et \texttt{for} interprète cela comme «~terminé~». Ce silence fait partie du protocole.}
\end{solution}

\begin{solution}{iq:iv3:python-and-pandas:8}
Il affiche \texttt{[-2147483648]}, \texttt{[255]} et \texttt{2147483648}. Les tableaux NumPy ont des types à largeur fixe et bouclent sans avertissement, alors que les entiers propres à Python grandissent au besoin.

Un tableau \texttt{int32} s'arrête à $2^{31} - 1$, et ajouter 1 le fait passer à la plus petite valeur. Un tableau \texttt{uint8} ne peut pas contenir de valeur négative, si bien que $1 - 2$ passe à 255. L'arithmétique sur les tableaux ne donne ni avertissement ni exception, car un contrôle par élément coûterait trop cher. Un scalaire NumPy se comporte autrement~: \texttt{np.int32(2**31 - 1) + 1} produit bien une \texttt{RuntimeWarning}, si bien qu'un contrôle sur une valeur isolée peut laisser croire que tout va bien.

Le problème se manifeste lorsque le type d'une colonne a été choisi pour une certaine taille, puis cumulé. Un volume cumulé en \texttt{int32} passe à un grand nombre négatif dès qu'il dépasse environ 2.1~milliards. Les horodatages en nanosecondes exigent 64 bits. La différence de numéros de séquence non signés donne une très grande valeur positive à la place d'une petite valeur négative.

\iqwrong{«~NumPy lèvera une erreur, comme le ferait Python.~» Python n'a aucun dépassement d'entier, donc il n'existe rien d'analogue à lever.}
\iqfollow{Comment vous protégez-vous~?}{On utilise \texttt{int64} pour tout ce qui est sommé, et l'on fait des assertions sur les totaux par un contrôle indépendant~: on compare la dernière valeur à un total calculé avec des entiers Python, ou l'on vérifie qu'un volume cumulé ne diminue jamais. La première différence de la somme cumulée ne prouve rien, car le débordement la ramène à la donnée d'entrée.}
\end{solution}

\begin{solution}{iq:iv3:python-and-pandas:9}
Seule C fonctionne, et elle renvoie la ligne dont \texttt{qty} vaut 2. A et B lèvent toutes deux une \texttt{ValueError}~: «~The truth value of a Series is ambiguous~».

Le \texttt{and} de Python doit réduire son membre de gauche à un seul vrai ou faux~; il appelle donc \texttt{bool} sur une Series, et pandas refuse, car une réponse par ligne ne peut pas former une seule valeur de vérité. Aucune bibliothèque ne peut redéfinir \texttt{and}. Pandas utilise donc l'opérateur bit à bit \texttt{\&}, qu'il peut surcharger pour travailler ligne par ligne.

B échoue à cause de la priorité. \texttt{\&} lie plus fortement que \texttt{>} et \texttt{<}, si bien que la ligne signifie \texttt{df.qty > (1 \& df.px) < 6}. C'est une comparaison enchaînée, que Python développe avec un \texttt{and}, donc la même erreur apparaît. Si les colonnes contenaient des flottants, \texttt{1 \& df.px} lèverait d'abord une \texttt{TypeError}.

C comporte les parenthèses~; chaque comparaison donne donc une Series de booléens, et \texttt{\&} les combine. La ligne dont \texttt{qty} vaut 2 a \texttt{px} égal à 5, ce qui passe. La ligne dont \texttt{qty} vaut 3 échoue sur \texttt{px}.

\iqfollow{Et avec \texttt{query}~?}{\texttt{df.query(\textquotesingle qty > 1 and px < 6\textquotesingle)} analyse lui-même le texte, si bien que \texttt{and} y fonctionne.}
\iqfollow{Et en NumPy pur~?}{C'est la même règle, avec un message semblable, dès qu'un tableau compte plus d'un élément~: \texttt{np.logical\_and} ou \texttt{\&} avec des parenthèses.}
\end{solution}

\begin{solution}{iq:iv3:python-and-pandas:10}
On n'y gagne rien, et souvent on perd un peu de vitesse. CPython 3.10 a un \index{verrou global de l'interpréteur}verrou global de l'interpréteur (GIL)~: un seul thread exécute du bytecode Python à la fois. Le verrou ne rend pas \texttt{count += 1} sûr, car cette instruction comporte plusieurs étapes de bytecode.

Le désassemblage de \texttt{count += 1} montre un chargement, une addition et un stockage, et rien dans le langage n'exclut qu'un changement de contexte survienne entre ces étapes. Quatre threads exécutant l'instruction nue n'ont perdu aucun des 8~millions d'incréments, sur trois essais ici avec CPython 3.10.12~: c'est le piège, puisque le test réussit, et cela ne garantit rien. Si l'on fait attendre les threads entre la lecture et l'écriture, notre test perd trois mises à jour sur quatre, à chaque fois. Il faut donc utiliser un \texttt{threading.Lock}, ou donner à chaque thread son propre compteur et les additionner à la fin.

Les threads aident là où le verrou est relâché~: en attente sur un socket ou un disque, et dans les opérations NumPy et de nombreuses extensions C, qui s'exécutent en parallèle sur de grands tableaux. Pour le calcul Python pur, on utilise des processus, au prix d'un envoi de données entre eux.

\iqwrong{«~Le GIL rend Python thread-safe.~» Il protège les structures propres à l'interpréteur, un bytecode à la fois. Une action composée dans le code n'est pas atomique.}
\iqfollow{Comment les processus partagent-ils un grand tableau~?}{Pas en le sérialisant avec pickle. Le module standard \texttt{shared\_memory} fournit un bloc nommé, et \texttt{np.ndarray} enveloppe son tampon dans chaque processus, si bien qu'ils voient tous une seule copie.}
\iqfollow{Est-ce définitif~?}{Python 3.13 a ajouté une variante optionnelle, dite free-threaded, sans verrou~; il faut donc vérifier quelle version est exécutée avant de se fier à l'un ou l'autre comportement.}
\end{solution}

\begin{solution}{iq:iv3:python-and-pandas:11}
Le \Cref{lst:iv3:pp:memo} fait passer le nombre d'appels d'une fonction de Fibonacci récursive pour $n = 30$ de 2\,692\,537 à 31. Un argument de type liste lève \texttt{TypeError}, car elle ne peut pas être hachée. \texttt{f(1, 2)} et \texttt{f(1, b=2)} sont deux entrées~: des réponses justes, calculées deux fois.
\omcode{code/interviews-3/02-python-and-pandas/python/iv_memo.py}{1}{14}{Un décorateur de mémoïsation.}\label{lst:iv3:pp:memo}

Sans cache, le nombre d'appels vérifie $c(n) = c(n - 1) + c(n - 2) + 1$, ce qui donne $2F_{n+1} - 1$, soit 2\,692\,537 pour $n = 30$. Avec le cache, chacune des 31 valeurs de 0 à 30 n'est calculée qu'une fois. Les appels récursifs passent par le nom décoré, et c'est ce qui fait fonctionner le mécanisme.

Le test \texttt{key not in cache} compte, car un résultat peut valoir 0, \texttt{None} ou une liste vide, et la forme \texttt{cache.get(key) or f(\dots)} recalcule dans chacun de ces cas, y compris pour Fibonacci de 0.

\texttt{wraps} recopie le nom et la docstring, afin qu'une trace de pile et un profilage montrent la vraie fonction, et non \texttt{inner}. Pour accepter des listes, on les convertit en tuples dans la clé, ou la mise en cache est laissée à l'appelant.

\iqwrong{Utiliser \texttt{id(arg)} dans la clé pour contourner le hachage. Une fois l'objet libéré, son adresse peut être réattribuée à une autre liste, et on lit alors la réponse mise en cache pour un autre appel.}
\iqfollow{Qu'utiliseriez-vous en production~?}{\texttt{functools.lru\_cache} avec une taille maximale~: ce cache croît indéfiniment, et sur une méthode il garde aussi en vie chaque \texttt{self}.}
\end{solution}

\begin{solution}{iq:iv3:python-and-pandas:12}
Le \Cref{lst:iv3:pp:stream} lit le fichier ligne par ligne et tient un total cumulé par symbole, si bien que la mémoire dépend du nombre de symboles distincts, et non de la taille du fichier.
\enlargethispage{-3\baselineskip}%
\omcode{code/interviews-3/02-python-and-pandas/python/iv_stream_csv.py}{1}{14}{Une ligne à la fois.}\label{lst:iv3:pp:stream}

\texttt{rows} est un générateur~: rien n'est lu tant que \texttt{notional\_by\_symbol} ne le demande pas, et chaque ligne peut être libérée dès qu'elle a été ajoutée. Dans notre test, sur un fichier de 200\,000 lignes, le pic d'allocation reste inférieur à 1~Mo, alors que \texttt{list(rows(path))} en utilise plus de 10~Mo, et cet écart croît avec la taille du fichier. Les prix restent exprimés en ticks entiers, si bien que les totaux sont exacts.

Si l'appelant s'arrête à mi-parcours, le fichier reste ouvert tant que le générateur est vivant. Il se ferme quand le générateur est fermé, par \texttt{close()} ou par le comptage de références de CPython lorsque la dernière référence disparaît. Les deux lèvent \texttt{GeneratorExit} au niveau du \texttt{yield}, ce qui exécute le \texttt{with}. D'autres implémentations de Python peuvent retarder cela, si bien qu'un appelant de longue durée doit appeler \texttt{close()} ou l'envelopper dans \texttt{contextlib.closing}.

\iqwrong{\texttt{pd.read\_csv(path)} sur une machine plus grosse. Cela construit toute la table en mémoire, ce que l'énoncé exclut. Avec \texttt{chunksize}, pandas peut le faire aussi, à condition d'additionner des sommes et de ne pas moyenner des moyennes.}
\iqfollow{Comment faites-vous pour qu'il résiste à une ligne défectueuse~?}{Compter les lignes défectueuses et poursuivre, jusqu'à une limite, afin qu'une seule ligne corrompue ne tue pas un traitement de 50~Go, et qu'un fichier corrompu ne passe pas sous silence.}
\end{solution}

\begin{solution}{iq:iv3:python-and-pandas:13}
Il affiche \texttt{[1]}, puis \texttt{[9] [1]}. \texttt{orders} est une seule liste de la classe, partagée par toutes les instances, si bien que \texttt{a.add(1)} l'ajoute à cette liste et que \texttt{b} la voit. L'affectation \texttt{a.orders = [9]} crée un nouvel attribut sur \texttt{a}, qui masque la liste de la classe pour \texttt{a} seul.

Ainsi, la lecture et l'ajout passent par l'objet de la classe, tandis que l'affectation écrit toujours sur l'instance. C'est toute l'incohérence. Pour un simple attribut de classe, la lecture cherche d'abord dans l'instance puis dans la classe, et l'affectation écrit dans l'instance. Une propriété, ou tout autre descripteur de données de la classe, prime sur les deux. La correction consiste à poser \texttt{self.orders = []} dans \texttt{\_\_init\_\_}.

Avec \texttt{\_\_slots\_\_ = (\textquotesingle name\textquotesingle, \textquotesingle orders\textquotesingle)}, les instances n'ont pas de \texttt{\_\_dict\_\_}, seulement des champs fixes. Dans notre test, une instance à slots occupe moins de la moitié de la mémoire d'une instance ordinaire, sous CPython 3.10. Une faute de frappe comme \texttt{self.qyt = 1} lève aussitôt \texttt{AttributeError}.

Un nom présent dans \texttt{\_\_slots\_\_} ne peut pas être aussi un attribut de classe avec une valeur par défaut~: Python lève \texttt{ValueError} lors de la construction de la classe, si bien que le bogue d'origine ne peut pas être reporté.

\iqwrong{«~\texttt{a.orders = [9]} modifie la liste pour tout le monde.~» Seule une modification de l'objet le fait. Cette réaffectation ne touche qu'une instance.}
\iqfollow{Et avec un million d'instances, faut-il vraiment utiliser une classe~?}{Souvent non. Les colonnes d'un tableau NumPy contiennent les mêmes données dans une fraction de la place.}
\end{solution}

\begin{solution}{iq:iv3:python-and-pandas:14}
Le \Cref{lst:iv3:pp:context} montre les deux formes, et chacune journalise «~release~» lorsque le corps lève une exception. Si \texttt{\_\_exit\_\_} renvoie une valeur vraie, l'exception est avalée, et l'exécution se poursuit après le bloc \texttt{with} comme si de rien n'était.
\omcode{code/interviews-3/02-python-and-pandas/python/iv_context.py}{1}{23}{Deux façons d'écrire un gestionnaire de contexte.\protect\hfil\protect\break}\label{lst:iv3:pp:context}

Python appelle \texttt{\_\_exit\_\_} quelle que soit la façon dont le bloc se termine, en lui passant les détails de l'exception s'il y en a une. Ne rien renvoyer, c'est-à-dire \texttt{None}, laisse l'exception se propager, de même que \texttt{False}. La règle sûre consiste donc à ne renvoyer une valeur vraie que lorsque ce gestionnaire entend traiter cette exception-là précisément, et à en tester le type.

Dans la version \texttt{contextlib}, une exception survenant dans le corps est levée à l'intérieur du générateur, au niveau du \texttt{yield}. Sans \texttt{try} ni \texttt{finally}, elle quitte le générateur à cet endroit, et la ligne suivant le \texttt{yield} ne s'exécute jamais. Un verrou reste donc pris, ou un fichier reste ouvert, dans le cas précis d'échec pour lequel le gestionnaire a été écrit. Le chemin normal ne montre rien, si bien qu'un test doit provoquer volontairement une erreur pour la détecter.

\iqwrong{Placer la libération dans \texttt{\_\_del\_\_}. Cette méthode s'exécute quand l'objet est collecté, si tant est qu'elle s'exécute, et à un moment qu'on ne peut prévoir.}
\iqfollow{Et si la libération lève elle-même une exception~?}{La nouvelle exception remplace l'ancienne, celle-ci étant rattachée comme son contexte. Il faut donc garder la libération simple, et journaliser l'exception d'origine.}
\end{solution}

\begin{solution}{iq:iv3:python-and-pandas:15}
Un champ \texttt{tags: list = []} lève \texttt{ValueError} à la définition de la classe, et non lorsqu'un ordre est créé. Le \Cref{lst:iv3:pp:dataclass} le corrige avec \texttt{field(default\_factory=list)}. Les objets d'une dataclass ordinaire ne peuvent pas entrer dans un ensemble, car elle n'a pas de hachage~: \texttt{@dataclass} met \texttt{\_\_hash\_\_} à \texttt{None}.
\omcode{code/interviews-3/02-python-and-pandas/python/iv_dataclass.py}{1}{14}{Une valeur par défaut mutable et une clé figée.\protect\hfil\protect\break}\label{lst:iv3:pp:dataclass}

L'erreur à la définition est la protection de Python contre la valeur par défaut partagée d'une fonction ordinaire~: sans elle, tous les ordres partageraient une seule liste. Une fabrique s'exécute une fois par instance.

Le hachage disparaît parce que \texttt{@dataclass} écrit un \texttt{\_\_eq\_\_} par valeur. Deux objets égaux doivent avoir le même hachage, et une valeur mutable pourrait changer après avoir été placée dans un ensemble, la laissant au mauvais endroit. Ainsi, la règle de Python est la suivante~: définir l'égalité sans déclarer l'objet immuable, c'est renoncer au hachage.

\texttt{frozen=True} fait lever une erreur à toute affectation d'un champ, et génère \texttt{\_\_hash\_\_} à partir des champs, si bien que \texttt{Key} fonctionne dans un ensemble ou comme clé de dictionnaire. Cela ne protège pas le contenu des champs. Un objet figé dont un champ est une liste peut encore recevoir des ajouts à cette liste, et le hacher lève alors \texttt{TypeError}, comme le fait hacher la liste elle-même.

\iqwrong{«~Ajouter \texttt{unsafe\_hash=True}.~» Cela s'exécute, et laisse un objet mutable se trouver dans un ensemble, où un changement ultérieur le fait perdre.}
\iqfollow{Pourquoi voudriez-vous \texttt{slots=True}~?}{À partir de 3.10, la classe est construite avec \texttt{\_\_slots\_\_}, ce qui réduit la mémoire par objet et supprime les fautes de frappe sur les attributs.}
\end{solution}

\begin{solution}{iq:iv3:python-and-pandas:16}
On trie deux fois, la clé la moins importante en premier, comme dans le \cref{lst:iv3:pp:sort}~: d'abord par notionnel décroissant, puis par symbole. Le tri de Python est stable, si bien que le second tri conserve l'ordre du premier à l'intérieur de chaque symbole.
\omcode{code/interviews-3/02-python-and-pandas/python/iv_sort.py}{1}{8}{Deux tris stables, la clé la moins importante en premier.}\label{lst:iv3:pp:sort}

Un tri stable n'échange jamais deux éléments de clés égales. Ainsi, après un tri par notionnel puis par symbole, les transactions d'un même symbole restent ordonnées par notionnel. \texttt{reverse=True} conserve aussi l'ordre des ex æquo, par définition du langage, si bien que ce premier passage n'inverse aucun élément à égalité. Écrire \texttt{sorted(trades, key=\dots)[::-1]} n'a pas cette propriété, car cela inverse aussi les ex æquo.

Lorsque la seconde clé est un nombre, un seul passage suffit~: \texttt{key=lambda t: (t[0], -t[1])}. Cela échoue pour une chaîne ou un \texttt{None}, dont aucun ne peut être changé de signe. La forme à deux passages fonctionne pour toute clé ordonnable, si bien qu'une colonne mêlant \texttt{None} et des nombres a toujours besoin d'une clé qui place le \texttt{None}.

Chaque passage coûte au pire $O(n \log n)$.

\iqwrong{Trier d'abord par symbole, puis par notionnel. Le dernier tri est la clé primaire, si bien qu'il ordonne toutes les transactions par notionnel et éparpille les symboles.}
\iqfollow{Et s'il y a plusieurs clés~?}{On les prend de la dernière à la première, ou on construit une seule fonction \texttt{key} renvoyant un tuple, dont chaque partie se trie comme on le souhaite.}
\end{solution}

\begin{solution}{iq:iv3:python-and-pandas:17}
Le premier affichage donne \texttt{[[-0.5, -2.5], [1.5, 1.5]]} et le second \texttt{[[-0.5, 0.5], [-1.5, 1.5]]}. Le second est juste. Chaque ligne doit totaliser zéro une fois la moyenne retirée, alors que dans le premier les lignes totalisent $-3$ et $3$.

\texttt{r.mean(axis=1)} a pour forme \texttt{(2,)}. La \index{diffusion}diffusion de NumPy aligne les formes à partir de la droite, si bien que \texttt{(2,)} correspond aux colonnes de \texttt{(2, 2)}, et elle retranche la moyenne du premier jour à la première action pour tous les jours, et la moyenne du second jour à la seconde action. C'est un résultat dénué de sens, sans erreur, parce que la matrice se trouve être carrée.

\texttt{keepdims=True} donne la forme \texttt{(2, 1)}, qui s'aligne sur les lignes et se répète sur les colonnes.

Avec trois actions, les formes sont \texttt{(2, 3)} et \texttt{(2,)}~: les dernières dimensions valent 3 et 2, qui ne correspondent pas, si bien que NumPy lève \texttt{ValueError}. Le bogue se cache donc lorsque le nombre de jours égale celui des actions, ce qui peut arriver avec une petite matrice de test ou une fenêtre courte, ou dans une fenêtre quelconque d'un calcul glissant. Il se cache aussi avec une seule action, où \texttt{(n, 1)} et \texttt{(n,)} se diffusent vers \texttt{(n, n)}.

\iqwrong{«~Cela lève une erreur si les formes ne correspondent pas.~» Seulement si elles ne correspondent pas après l'alignement par la droite, et cette règle laisse précisément passer une matrice carrée.}
\iqfollow{Comment vous protégez-vous~?}{Garder \texttt{keepdims=True} comme réflexe, et tester avec des entrées non carrées de plus d'une action, afin qu'un mauvais axe échoue bruyamment. Une vérification selon laquelle les lignes totalisent zéro le détecte sur toute forme.}
\end{solution}

\begin{solution}{iq:iv3:python-and-pandas:18}
Il affiche \texttt{[ 0 99  2  3  4]}. Seul \texttt{b} est une vue~: une tranche partage la mémoire avec \texttt{a}, alors que l'indexation par une liste ou un masque booléen crée une copie, si bien que les écritures via \texttt{c} et \texttt{d} n'atteignent pas \texttt{a}.

La règle découle de ce qui peut être exprimé. Une tranche est un début, une fin et un pas, si bien qu'un nouvel en-tête au-dessus du même tampon peut la décrire, pour un coût en $O(1)$. Des positions quelconques ne peuvent pas être décrites ainsi, si bien que NumPy les copie. \texttt{np.shares\_memory(a, b)} indique dans quel cas on se trouve.

\texttt{small = huge[:10]} garde les 10~Go en mémoire, puisque \texttt{small.base} pointe vers le tampon d'origine. On utilise \texttt{huge[:10].copy()}, puis on laisse \texttt{huge} disparaître.

\texttt{df[df.x > 0][\textquotesingle y\textquotesingle] = 1} avec pandas~2.3 laisse \texttt{df} inchangé, avec un \texttt{SettingWithCopyWarning}. Le masque donne un nouvel objet, et l'affectation porte sur cet objet. Celui-ci est ensuite jeté. On écrit \texttt{df.loc[df.x > 0, \textquotesingle y\textquotesingle] = 1}, qui est une seule opération sur \texttt{df}.

\iqwrong{«~Pandas copie toujours.~» Ce n'est pas fiable dans un sens ou dans l'autre, et c'est pourquoi l'avertissement existe. Pandas~3 fait du copy-on-write le comportement par défaut, donc il faut vérifier la version utilisée.}
\iqfollow{Quand une vue est-elle un bogue et non un gain~?}{Quand une fonction modifie un tableau qu'on lui a transmis~: la modification atteint les données de l'appelant. Dans le code partagé, on rend les entrées en lecture seule, avec \texttt{arr.flags.writeable = False}.}
\end{solution}

\begin{solution}{iq:iv3:python-and-pandas:19}
(A, 1) donne 10.75, (B, 1) donne 20 et (A, 2) donne NaN, comme dans le \cref{lst:iv3:pp:vwap}. \texttt{groupby(\dots).price.mean()} donne 10.5 pour (A, 1), ce qui traite de la même façon une transaction de 100~actions et une transaction de 300~actions.
\omcode{code/interviews-3/02-python-and-pandas/python/iv_vwap.py}{1}{7}{Deux sommes par groupe et une seule division.}\label{lst:iv3:pp:vwap}

La moyenne pondérée vaut $(10 \times 100 + 11 \times 300) / 400 = 4300 / 400 = 10.75$. On calcule les deux sommes par groupe, puis on divise. Les deux résultats ont le même indice de groupe, si bien que la division s'aligne par groupe sans travail supplémentaire.

Le NaN pour (A, 2) est la réponse honnête~: aucune action n'a été échangée, donc aucune moyenne n'existe, et 12 serait un prix auquel rien n'a été fait. Pandas donne NaN pour $0 / 0$ ici, et il n'y a pas besoin de garde. Une version Python pure lèverait une exception, et une version C++ diviserait par zéro.

L'alternative habituelle est \texttt{groupby(\dots).apply(lambda g: \dots)}, qui est correcte mais exécute une fonction Python une fois par groupe, alors que les deux sommes du listing s'exécutent en code compilé.

\iqwrong{\texttt{df.price * df.qty} sommé puis divisé par la quantité de toute la table. Cela donne un seul nombre pour la table, et non un nombre par groupe.}
\iqfollow{Comment le replacer sur chaque ligne~?}{On utilise \texttt{transform(\textquotesingle sum\textquotesingle)} à la place de \texttt{sum}, ce qui conserve l'index d'origine.}
\end{solution}

\begin{solution}{iq:iv3:python-and-pandas:20}
Par défaut, la transaction à 1 n'obtient aucune cotation, celle à 5 obtient la cotation à 5, et celle à 10 obtient la cotation à 9. Le \cref{lst:iv3:pp:asof} est l'appel usuel. Avec \texttt{allow\_exact\_matches=False}, la transaction à 5 obtient la cotation à 2. Avec \texttt{direction=\textquotesingle forward\textquotesingle}, elles obtiennent 2, 5 et aucune cotation, ce qui constitue le bogue du backtest.
\omcode{code/interviews-3/02-python-and-pandas/python/iv_merge_asof.py}{1}{8}{La dernière cotation à l'instant de chaque transaction ou avant.\protect\hfil\protect\break}\label{lst:iv3:pp:asof}

Un appariement vers l'avant donne à chaque transaction une cotation issue de son futur. Un signal construit sur cette base paraît excellent, mais il ne peut pas rapporter, puisqu'à l'instant~1 personne ne connaît encore la cotation de l'instant~2.

Une correspondance exacte est un choix sur le sens des horodatages. Si une cotation à 5 est publiée avant une transaction à 5, on la garde. Si les horloges ne permettent pas de savoir laquelle est venue en premier, alors \texttt{allow\_exact\_matches=False} est le choix prudent, et il coûte peu.

Les deux entrées doivent être triées selon la clé, et pandas lève une \texttt{ValueError} («~left keys must be sorted~») si elles ne le sont pas, ce qui explique que le listing trie d'abord. Avec \texttt{by=\textquotesingle symbol\textquotesingle}, les appariements restent à l'intérieur d'un symbole.

En pratique, on ajoute \texttt{tolerance}. Sinon, une transaction survenant après un long écart reçoit une cotation vieille de plusieurs heures, et rien ne le signale.

\iqwrong{\texttt{merge} sur la colonne temporelle, puis remplissage vers l'avant. Cela ne fonctionne que si l'on trie et groupe correctement, cela remplit à travers les symboles si l'on ne groupe pas, et cela coûte bien plus de mémoire.}
\iqfollow{Quel est le coût~?}{Après le tri, un seul parcours des deux~: $O(m + n)$ pour la fusion, et $O(m \log m + n \log n)$ si elles arrivent non triées.}
\end{solution}

\begin{solution}{iq:iv3:python-and-pandas:21}
La barre de 09:30 regroupe les ticks allant de 09:30:00 jusqu'à 09:31:00 non compris~: ouverture 10, plus haut 12, plus bas 10, clôture 12. Par défaut, pandas étiquette cette barre 09:30, le début de la barre. Une jointure sur cette étiquette traite 12, un prix observé à 09:30:50, comme connu à 09:30:00. Le \cref{lst:iv3:pp:bars} étiquette chaque barre à sa fin.
\omcode{code/interviews-3/02-python-and-pandas/python/iv_bars.py}{1}{7}{Barres étiquetées lorsqu'elles sont complètes.}\label{lst:iv3:pp:bars}

Une barre n'est donc connue qu'une fois clôturée, et \texttt{label=\textquotesingle right\textquotesingle} l'exprime~: la même barre devient 09:31. Une stratégie qui lit à 09:31 la barre étiquetée 09:31 ne connaît alors que le passé. \texttt{closed} est une question distincte, celle de savoir de quel côté de la minute tombe un tick survenant exactement à 09:31:00, et ici il reste à \texttt{left}.

La barre de 09:32 ne contient aucun tick, si bien que les quatre colonnes valent NaN. La clôture d'une barre vide peut raisonnablement reprendre le dernier prix, comme le fait le listing, puisque c'est encore le dernier prix connu, ce qui donne 11. L'ouverture, le plus haut et le plus bas doivent rester NaN ou être égaux à ce prix, et toute colonne de volume doit valoir 0. Une stratégie doit aussi savoir que la barre était vide, car un prix rempli n'est pas un marché négociable à ce moment-là.

\iqwrong{Remplir toutes les colonnes vers l'avant sans les marquer. Une barre remplie ressemble à un marché calme, et une estimation de la volatilité sort alors trop faible.}
\iqfollow{Qu'en est-il d'une frontière de séance~?}{\texttt{resample} produit aussi des barres pendant la nuit, donc on filtre sur les heures de séance, ou on groupe par jour, avant de remplir.}
\end{solution}

\begin{solution}{iq:iv3:python-and-pandas:22}
Il affiche \texttt{[nan, 12.0, 23.0, nan]}, \texttt{[11, 22, 33]} et \texttt{[1.0, 12.0, 23.0, 30.0]}. La deuxième ligne est le bogue qui ne lève aucune erreur. La première est le comportement voulu de pandas~: il aligne par index, sur l'union des deux.

\texttt{a + b} apparie les étiquettes~: à 2 il additionne 2 et 10, à 3 il additionne 3 et 20, et les étiquettes présentes d'un seul côté n'ont pas de partenaire, d'où NaN. \texttt{a.values + b.values} jette les étiquettes et additionne par position, si bien que l'étiquette~1 de l'une est additionnée à l'étiquette~2 de l'autre.

Pour une paire, cela revient à comparer le prix du lundi d'une action au prix du mardi de l'autre, à partir du jour manquant et pour chaque jour suivant. Les spreads restent plausibles, la corrélation baisse, ou un décalage apparent de type avance-retard apparaît. Il n'y a ni NaN ni avertissement, et les formes concordent, car la seconde série dispose d'un jour de plus pour compléter le compte.

Il ne faut donc jamais quitter l'index lorsqu'il y a deux séries. On réindexe les deux sur un même calendrier, on décide ce que signifie un jour manquant, et l'on conserve le NaN tant qu'on n'a pas décidé. \texttt{add(b, fill\_value=0)} convient aux quantités, où 0 signifie aucune. Pour des prix, cela produit une valeur absurde, et un mouvement de 100~\%.

\iqwrong{«~Le NaN signifie que pandas est cassé, il suffit d'utiliser \texttt{.values}.~» Le NaN est l'avertissement que vous cherchiez à éviter.}
\iqfollow{Et comment le voir dans une grande table~?}{On vérifie par une assertion que les deux index sont égaux avant toute étape positionnelle, ou on contrôle \texttt{a.index.difference(b.index)}.}
\end{solution}

\begin{solution}{iq:iv3:python-and-pandas:23}
On obtient un ratio de Sharpe annualisé d'environ 21, sur du bruit pur. La correction est \texttt{shift}~: \texttt{signal.shift(1) * ret}, qui donne environ 0. Le rendement du jour figurait dans le signal, si bien que le backtest trade sur le journal de demain. Le \cref{lst:iv3:pp:noise} l'exécute.
\omcode{code/interviews-3/02-python-and-pandas/python/iv_noise_backtest.py}{1}{11}{Un backtest sur du bruit, avant et après le décalage.\protect\hfil\protect\break}\label{lst:iv3:pp:noise}

Avec $r \sim N(0, \sigma^2)$, le profit vaut $\text{sign}(r) \times r = |r|$. Sa moyenne est $\sigma\sqrt{2/\pi}$, et sa variance vaut $\sigma^2(1 - 2/\pi)$. Le ratio quotidien vaut $\sqrt{2/\pi} / \sqrt{1 - 2/\pi} = 1.3236$, et multiplier par $\sqrt{252}$ donne 21.0. Le résultat est le même quel que soit $\sigma$, puisqu'un ratio de Sharpe est sans unité.

Notre simulation sur 25\,200 jours donne 21.1 et, avec le décalage, $-0.03$. Du bruit passé dans la version décalée ne devrait rien montrer, et c'est ce test qu'il faut conserver.

Une bonne habitude pour trouver les bogues en découle. On exécute chaque stratégie sur du bruit mélangé ou simulé. Tout résultat nettement supérieur à zéro sur ce bruit signale une fuite de données.

Pour rendre l'erreur difficile à écrire, on donne à la bibliothèque de recherche un seul endroit qui joint les signaux et les rendements, et qui applique lui-même le décalage, de sorte qu'une stratégie ne puisse renvoyer que des positions décidées avec les données allant jusqu'à chaque date. Un biais d'anticipation devra alors être écrit délibérément.

\iqwrong{«~Un Sharpe de 21 signifie une excellente stratégie, il faut donc vérifier les coûts.~» Tout ratio de Sharpe annuel très supérieur à 3 ou 4 sur des données quotidiennes doit d'abord faire chercher un bogue.}
\end{solution}

\begin{solution}{iq:iv3:python-and-pandas:24}
Il affiche \texttt{4.0 2.0 2 3 1.33}, puis \texttt{0.0 nan}, puis \texttt{\{\textquotesingle A\textquotesingle: 4\}}, puis \texttt{False}.

\texttt{sum}, \texttt{mean} et \texttt{count} ignorent tous les NaN par défaut, si bien que la moyenne vaut $4 / 2$, et \texttt{count} vaut 2 tandis que \texttt{len} vaut 3. \texttt{fillna(0)} modifie d'abord le sens du calcul~: la moyenne vaut désormais $4 / 3$. La somme des seuls NaN vaut 0, par ce même mécanisme d'exclusion, ce qui peut faire passer un jour manquant pour un vrai zéro~; \texttt{min\_count=1} la rend NaN.

\texttt{groupby} écarte par défaut les lignes dont la clé est manquante, si bien que la ligne de valeur 2 disparaît du total, sans aucun avertissement. On utilise \texttt{dropna=False} pour la conserver. Et \texttt{s == np.nan} est faux partout, donc on utilise \texttt{s.isna()}.

\texttt{prices.fillna(0)} affirme que le prix est tombé à zéro. Le rendement de ce jour vaut $-100\%$, et le jour suivant est une division par zéro, que pandas donne sous la forme \texttt{inf}. Les sommes et les moyennes de ces colonnes deviennent alors inutilisables. Pour les prix, on laisse le NaN, ou on reporte le dernier prix vers l'avant avec \texttt{ffill} si un prix périmé est acceptable, et on le dit. On calcule les rendements après cette étape, et on ne remplit pas les rendements par la moyenne, car cela invente des données.

\iqwrong{«~Un NaN vaut zéro dans une somme, donc le remplir est sans danger.~» C'est vrai dans une somme, mais pas dans une moyenne, un comptage, un produit ou une division.}
\end{solution}

\begin{solution}{iq:iv3:python-and-pandas:25}
On utilise \texttt{category} pour \texttt{symbol}, on garde \texttt{price} en \texttt{float64}, et \texttt{qty} en \texttt{int64} ou \texttt{int32} selon son étendue. Les prix occupent 800~Mo, les quantités 800~Mo ou 400~Mo, et les symboles 200~Mo, contre environ 1~Go pour une colonne de chaînes ordinaire issue de \texttt{read\_csv} et environ 6~Go avec une chaîne par ligne.

Une colonne de chaînes ordinaire contient 8~octets par ligne pour le pointeur, plus les objets chaînes. Construite à partir du module \texttt{csv} de Python, chaque ligne obtient sa propre chaîne, d'environ 53~octets pour une chaîne courte sous CPython~3.10, ce qui fait environ 6~Go. L'analyseur C de \texttt{pd.read\_csv} partage les chaînes égales, donc on est plus proche de 1~Go, bien que \texttt{memory\_usage(deep=True)} annonce toujours environ 6~Go, car il compte chaque élément. \texttt{category} stocke chaque chaîne distincte une seule fois, et un petit code entier par ligne. Pour 4\,000 valeurs, le code est un \texttt{int16}~: 2~octets par ligne, soit 200~Mo. Dans notre test sur 200\,000 lignes avec une chaîne par ligne, c'est plus de dix fois plus petit, et face aux chaînes partagées de \texttt{read\_csv}, cela fait environ quatre fois d'économie.

Un \texttt{float32} a 24~bits de mantisse. Un prix de 101.27 se relit comme 101.26999664\dots, si bien qu'une égalité avec une grille de tick échoue, et un prix flottant voisin de 10\,000 a un pas d'environ 0.001. Pour les totaux, c'est bien pire. Au-delà de $2^{24} = 16\,777\,216$, tous les entiers ne sont pas représentables, donc ajouter 1 à cette valeur ne la change pas~: un volume cumulé cesse simplement de croître, sans aucune erreur.

\iqwrong{«~On réduit tout au plus petit type qui convient aujourd'hui.~» Les données du lendemain peuvent faire déborder un \texttt{int32} sans avertissement, car l'arithmétique sur les tableaux reboucle.}
\iqfollow{Et si les chaînes se répètent peu~?}{Alors \texttt{category} est moins bon~: il stocke toutes les chaînes et les codes en plus. Il convient aux colonnes à faible cardinalité.}
\end{solution}

\begin{solution}{iq:iv3:python-and-pandas:26}
Le \Cref{lst:iv3:pp:checks} est un tel fichier de tests, et il passe face à la fonction. \texttt{0.1 + 0.2} vaut 0.30000000000000004, car aucune des trois valeurs n'a de représentation binaire exacte. \texttt{pytest.approx} accepte une tolérance relative de $10^{-6}$ ou une tolérance absolue de $10^{-12}$, selon la plus grande des deux. Ainsi \texttt{1e-13 == approx(0)} passe, et \texttt{1e-9 == approx(0)} échoue.
\omcode{code/interviews-3/02-python-and-pandas/python/vwap_checks.py}{1}{23}{Un seul corps de test, plusieurs cas.}\label{lst:iv3:pp:checks}

\texttt{parametrize} exécute un seul corps sur une table de cas, et il rapporte chacun comme un test distinct, si bien qu'une ligne en échec ne masque pas les autres. Les cas sont choisis pour faire échouer un modèle faux~: 10.75 fait échouer la moyenne simple, qui donne 10.5.

Une tolérance relative est inutile près de zéro, car $10^{-6}$ de zéro vaut zéro, et c'est pourquoi la limite absolue existe. Face à une cible de 0, tout ce qui reste sous $10^{-12}$ passe, et une valeur comme $10^{-9}$ échoue même s'il ne s'agit que d'un bruit d'arrondi, donc on passe \texttt{abs=} lorsque l'échelle l'exige.

Il faut ajouter les cas qu'une table de lignes ordinaires laisse échapper~: quantité nulle, une seule ligne, une table vide, plusieurs groupes à la fois, et un groupe qui se présente en deux blocs séparés. Il faut aussi vérifier une limite dont la réponse est connue~: si tous les prix sont égaux, la moyenne vaut ce prix quelles que soient les tailles.

\iqwrong{\texttt{assert abs(a - b) < 1e-9} partout. Sur une valeur proche de $10^{12}$, cette limite est plus serrée que le type ne peut l'exprimer, si bien qu'un calcul correct échoue, et près de $10^{-12}$ elle laisse passer n'importe quoi.}
\end{solution}

\begin{solution}{iq:iv3:python-and-pandas:27}
Il affiche \texttt{2 4 2.67}, puis le double exact, \texttt{2.67499999999999982\dots}, puis \texttt{False}, et enfin \texttt{2.68}. La fonction \texttt{round} de Python envoie les moitiés vers le voisin pair, et 2.675 sous forme de double n'est pas du tout une moitié.

L'arrondi au pair évite que de longues colonnes de valeurs arrondies ne dérivent vers le haut, si bien que \texttt{round(2.5)} vaut 2 et \texttt{round(3.5)} vaut 4. La fonction \texttt{round} de NumPy envoie elle aussi les moitiés exactes vers le voisin pair, si bien que «~elle arrondit vers le haut~» est faux dans les deux cas. Avec un nombre de décimales, elle commence par multiplier par une puissance de dix, si bien que sous NumPy 2.2 \texttt{np.round(2.675, 2)} vaut 2.68, là où Python donne 2.67.

\texttt{round(2.675, 2)} ne serait un cas d'égalité que si l'entrée valait vraiment 2.675. Le double le plus proche se trouve juste en dessous, comme le montre la deuxième ligne, si bien que l'arrondi correct de ce double donne 2.67. \texttt{Decimal(2.675)} convertit le double exactement et met le problème en évidence, ce qui explique que \texttt{Decimal(0.1)} ne soit pas égal à \texttt{Decimal(\textquotesingle 0.1\textquotesingle)}.

Il faut donc que les frais arrivent sous forme de texte, ou sous forme d'un nombre entier d'unités plus petites. \texttt{Decimal(\textquotesingle 2.675\textquotesingle)} contient la valeur exacte, et \texttt{quantize} avec \texttt{ROUND\_HALF\_UP} donne 2.68. C'est la quatrième ligne. Il faut énoncer la règle d'arrondi dans le code, puisque c'est un barème de frais ou un recueil de règles qui la fixe, et non le langage.

\iqwrong{\texttt{round(x + 0.005, 2)}, ou l'ajout d'un petit epsilon. Cela marche sur certaines valeurs et casse sur d'autres, et cela introduit un biais qui, sur les exécutions d'une journée, peut devenir de l'argent réel.}
\iqfollow{Et sur un chemin critique~?}{On y garde des entiers exprimés en unités plus petites, et on utilise \texttt{Decimal} aux bords, là où les montants sont déclarés ou facturés.}
\end{solution}

\begin{solution}{iq:iv3:python-and-pandas:28}
\texttt{sum} renvoie 0.0, et \texttt{math.fsum} renvoie 1.0, la vraie réponse. Les nombres en double précision proches de $10^{16}$ sont espacés de 2, si bien que $10^{16} + 1$ revient à $10^{16}$, et le 1 a disparu avant la soustraction. Le \Cref{lst:iv3:pp:neumaier} est la forme de Neumaier de la \index{sommation compensée}sommation compensée, qui renvoie aussi 1.0. La version classique de Kahan renvoie 0.0 sur cette liste.
\omcode{code/interviews-3/02-python-and-pandas/python/iv_exact_sum.py}{1}{10}{Conserver ce que l'addition arrondit.\protect\hfil\protect\break}\label{lst:iv3:pp:neumaier}

Chaque addition perd quelques bits de poids faible, et ces bits peuvent être calculés exactement~: pour deux nombres en double précision, la plus grande valeur moins la somme arrondie, plus la plus petite, est exactement ce qui a été perdu. Le listing ajoute ces termes à un total séparé et petit, et ne réunit les deux qu'à la fin.

Kahan reporte la correction sur l'entrée suivante, si bien que la valeur suivante, $-10^{16}$, avec sa correction de $+1$, se place à $-10^{16} + 1$, qui s'arrondit aussitôt. La perte est réelle~: sa correction suppose que le total courant est plus grand que chaque nouveau terme, et ici le nouveau terme est aussi grand que le total.

Aucune des deux n'est exacte dans tous les cas, alors que \texttt{fsum} l'est~: elle conserve assez de sommes partielles pour donner le résultat correctement arrondi, au prix d'une certaine perte de vitesse. Sur un million de copies de 0.1, la \texttt{sum} classique s'écarte d'environ $1.3 \times 10^{-6}$, tandis que \texttt{fsum} et le listing donnent tous deux 100\,000.0.

\iqwrong{«~Trier d'abord les nombres.~» Additionner du plus petit au plus grand aide pour des valeurs d'un même signe, mais cela ne change rien ici, où les deux grandes valeurs ne s'annulent qu'une fois la petite perdue.}
\iqfollow{Une version plus récente de Python modifie-t-elle cela~?}{Oui~: à partir de la version 3.12, la fonction \texttt{sum} intégrée sur des flottants utilise un algorithme plus précis, si bien que la première réponse peut changer selon la version.}
\end{solution}

\begin{solution}{iq:iv3:python-and-pandas:29}
Le heartbeat s'arrête pendant les 200 ms entières. Une boucle d'événements est coopérative~: sur un même thread, une coroutine s'exécute jusqu'à atteindre un \texttt{await} qui cède la main, et le travail bloquant n'en a aucun. Le \Cref{lst:iv3:pp:async} montre cela sous la forme d'une suite d'événements~: avec un traitement bloquant, tout le travail s'exécute avant le heartbeat suivant.
\omcode{code/interviews-3/02-python-and-pandas/python/iv_async.py}{1}{27}{Un gestionnaire bloquant et un gestionnaire qui cède la main.\protect\hfil\protect\break}\label{lst:iv3:pp:async}

Dans le journal du test, \texttt{handler\_blocks} produit un heartbeat, puis trois morceaux de travail, puis deux heartbeats. Aucune horloge n'est nécessaire pour le montrer. Une bourse qui ne voit aucun heartbeat pendant aussi longtemps peut couper la session, et toutes les autres coroutines sont elles aussi à l'arrêt, y compris celles qui lisent les données de marché.

Pour une entrée ou une sortie bloquante, il faut utiliser une bibliothèque asynchrone, ou \texttt{asyncio.to\_thread}, qui exécute l'appel sur un thread de travail et cède la main à la boucle pendant l'attente. Comme le verrou global de l'interpréteur (GIL) est relâché pendant une attente réelle, la boucle continue de tourner. Pour un calcul Python pur, cela n'aide pas longtemps à cause de ce verrou~: le thread reste en concurrence avec la boucle. Il faut donc le confier à un processus, avec \texttt{run\_in\_executor} et un \texttt{ProcessPoolExecutor}.

\iqwrong{«~Ajouter \texttt{await asyncio.sleep(0)} à la fin de la fonction.~» Cela cède la main une seule fois, après que le mal est fait. L'éparpiller dans le travail aide bien, car cela découpe le délai en morceaux plus petits, au prix d'un code plus encombré.}
\iqfollow{Et si vous appelez une coroutine sans \texttt{await}~?}{Cela crée un objet coroutine et n'exécute rien. Quand cet objet est récupéré, Python émet l'avertissement «~was never awaited~». L'ordre ne part tout simplement jamais.}
\end{solution}

\begin{solution}{iq:iv3:python-and-pandas:30}
La bonne pente vaut 2, calculée sur les trois lignes complètes, et le \cref{lst:iv3:pp:ols} la calcule. En centrant chaque colonne séparément, on obtient 0.00056. Un symbole à une seule ligne renvoie NaN, de même qu'un symbole dont \texttt{x} est constant.
\omcode{code/interviews-3/02-python-and-pandas/python/iv_group_ols.py}{1}{11}{Une pente par symbole, à partir des sommes par groupe.\protect\hfil\protect\break}\label{lst:iv3:pp:ols}

La pente est~:
\[ \sum (x - \bar x)(y - \bar y) \Big/ \sum (x - \bar x)^2, \]
où les deux moyennes sont prises sur les lignes utilisées. \texttt{transform(\textquotesingle mean\textquotesingle)} donne à chaque ligne la moyenne de son groupe, et le reste se réduit à de l'arithmétique sur les colonnes et à deux sommes par groupe, le tout en code compilé.

L'erreur vient de la ligne (100, NaN). Colonne par colonne, pandas prend $\bar x = 26.5$ sur les quatre lignes et ajoute le 100 au dénominateur, tandis que le numérateur écarte cette ligne parce que son produit vaut NaN. Les deux sommes portent donc sur des lignes différentes, et rien n'en avertit. Supprimer d'abord les lignes incomplètes fait porter chaque somme sur les mêmes lignes.

Une pente calculée à partir d'une seule ligne ou d'un \texttt{x} constant n'a aucun sens, et NaN le signale. Les deux cas donneraient $0 / 0$ en arithmétique exacte, mais en flottants la moyenne d'un groupe peut être décalée d'un bit, et un \texttt{x} constant donne alors une pente finie qui n'est que de l'arrondi. Le listing renvoie donc aussi NaN lorsque la somme des $dx^2$ est trop petite, par rapport à la somme des $x^2$, pour que le résultat ait un sens. Un filtre sur un nombre minimal de lignes doit venir après cette étape, car une pente issue de 3 lignes s'affiche exactement comme une pente issue de 3\,000.

\iqwrong{\texttt{groupby(\textquotesingle symbol\textquotesingle).apply(lambda g: np.polyfit(g.x, g.y, 1))}. C'est juste sur des données propres, mais cela exécute une fonction Python 5\,000 fois~; un NaN dans un groupe gâche l'ajustement de ce groupe.}
\iqfollow{Comment obtenez-vous l'ordonnée à l'origine~?}{\texttt{means.y - beta * means.x}, calculée par symbole sur les mêmes lignes complètes.}
\end{solution}

\begin{solution}{iq:iv3:python-and-pandas:31}
Retirer d'abord les lignes impossibles, puis trier, puis supprimer les doublons, puis tester les valeurs aberrantes par rapport aux seules valeurs passées, et rééchantillonner à l'intérieur de chaque symbole et de chaque journée, comme dans le \cref{lst:iv3:pp:ticks}. Avec une limite de 5~\% sur une médiane de 20, un vrai saut de 6~\% perd ses 10 premiers ticks, et le reste passe.
\enlargethispage{-3\baselineskip}%
\omcode{code/interviews-3/02-python-and-pandas/python/iv_clean_ticks.py}{1}{19}{Nettoyage des cotations, puis rendements à une minute.\protect\hfil\protect\break}\label{lst:iv3:pp:ticks}

L'ordre compte. Un côté nul laissé en place tirerait la médiane, et une ligne tardive laissée non triée serait comparée aux mauvais voisins. \texttt{keep=\textquotesingle last\textquotesingle} après un tri stable conserve la dernière arrivée à un même horodatage, ce qui est un choix qui mérite d'être énoncé.

Une fenêtre centrée compare un tick à des ticks arrivés après lui. Dans un backtest, cela élimine de mauvais prix enregistrés en utilisant une information que l'on n'avait pas à ce moment-là, et rend l'historique nettoyé plus lisse que tout ce qu'on pourrait réellement négocier.

La clé de jour empêche \texttt{diff} de franchir une nuit. Le remplissage vers l'avant est limité à 5 minutes, si bien qu'un mid périmé n'est pas présenté comme un marché calme pendant des heures.

Les cinq premiers ticks d'un symbole n'ont pas de médiane passée et passent sans contrôle, si bien qu'une mauvaise valeur y survit. Il faut les tester par rapport à un prix de référence, par exemple la clôture de la veille.

Lors du saut, la médiane passée reste au vieux niveau jusqu'à ce que 10 des 20 derniers mids soient nouveaux. Ces ticks ressemblent à des erreurs, si bien qu'un vrai mouvement et un mauvais prix enregistré commencent de la même manière. Il faut donc compter les rejets par symbole et émettre une alerte à ce sujet, et ne jamais laisser le filtre tourner en silence, car une longue série de rejets est une information en soi.

\iqwrong{Corriger la valeur aberrante en plafonnant les rendements à la fin. Le mauvais mid donne alors deux rendements plafonnés, l'un à la hausse et l'autre à la baisse, et ils ressemblent à de la volatilité.}
\end{solution}

\begin{solution}{iq:iv3:python-and-pandas:32}
Avec $a = 0.1$, le réglage par défaut donne 2.0701 à la troisième valeur, et \texttt{adjust=False} donne 1.29. Le \Cref{lst:iv3:pp:ewma} conserve deux nombres par flux et coïncide avec pandas dans les tests, y compris sur NaN. Les deux formes concordent dès que $(1 - a)^n$ est négligeable.
\omcode{code/interviews-3/02-python-and-pandas/python/iv_ewma.py}{4}{40}{Deux nombres d'état par flux.}\label{lst:iv3:pp:ewma}

Le réglage par défaut est une véritable moyenne pondérée, $\sum_i (1 - a)^i x_{n-i} / \sum_i (1 - a)^i$. Ici cela donne $(3 + 0.9 \times 2 + 0.81 \times 1) / 2.71 = 5.61 / 2.71$. Le numérateur et le dénominateur admettent chacun une mise à jour sur une seule ligne, et c'est là toute l'astuce. \texttt{adjust=False} est la récurrence $y_n = (1 - a) y_{n-1} + a x_n$ initialisée en $x_0$, si bien que la première valeur conserve un poids de $(1 - a)^n$ et non la faible part qu'elle aurait dans un long historique. Dans le test, les deux diffèrent de moins de $10^{-6}$ après 200 valeurs.

Pour un NaN, pandas fait vieillir les anciens poids sans rien ajouter, ce que reproduit la première classe. La moyenne elle-même ne bouge pas à cette étape, mais la valeur suivante compte davantage.

Pour des temps irréguliers, utiliser une décroissance de $e^{-\lambda \, \Delta t}$ avec $\lambda = \ln 2 / h$. La seconde classe coïncide avec \texttt{ewm} appelé avec \texttt{halflife} et \texttt{times}, dans le test, NaN compris. Aux extrêmes, elle se comporte de façon raisonnable. Un horodatage commun donne $\Delta t = 0$, si bien que les deux ticks reçoivent un poids égal. Après un écart bien plus long que la demi-vie, les anciens poids ont presque disparu, si bien que la moyenne saute vers le nouveau tick.

\iqwrong{La récurrence à $a$ fixe sur des ticks irréguliers. Une rafale de 100 ticks en une milliseconde efface alors le passé, ce qui fait dépendre la moyenne du débit des messages.}
\end{solution}

\begin{solution}{iq:iv3:python-and-pandas:33}
\texttt{z @ L} a pour covariance $L^\top L$, ce qui, pour 0.8, donne des variances de 1.64 et 0.36 avec une covariance de 0.48. La bonne ligne est \texttt{z @ L.T}, comme dans le \cref{lst:iv3:pp:corr}. La matrice avec 0.9, 0.9 et $-0.9$ fait lever à \texttt{cholesky} une \texttt{LinAlgError}, et une corrélation égale à 1 exactement produit le même effet.
\omcode{code/interviews-3/02-python-and-pandas/python/iv_correlated.py}{1}{15}{Normales corrélées, et une réparation.\protect\hfil\protect\break}\label{lst:iv3:pp:corr}

On prend une ligne $z$ de covariance $I$. Alors $x = z L^\top$ a pour covariance $\text{Cov}(x) = L\, I\, L^\top = C$. L'énoncé usuel $x = L z$ revient à la même chose, avec $z$ en colonne~; c'est de là que vient l'erreur. Pour 0.8, $L$ a pour lignes $(1, 0)$ et $(0.8, 0.6)$. Les deux versions donnent une corrélation positive, 0.8 contre 0.625, et la fausse version casse aussi les variances unitaires, si bien que tester la covariance empirique le détecte en une ligne.

Si A évolue étroitement avec B et avec C, alors B et C ne peuvent pas évoluer en sens opposés. Cela se voit sur les valeurs propres, qui sont ici $-0.8$, 1.9 et 1.9. La variance d'un portefeuille ne peut pas être négative, si bien qu'aucune donnée ne peut produire cette matrice si toutes les paires utilisent les mêmes dates. De telles matrices naissent de paires estimées sur des fenêtres différentes, ou de retouches à la main. Il faut donc d'abord regarder la source. Une réparation rapide écrête les valeurs propres négatives et renormalise la diagonale à 1, ce qui donne ici 0.5, 0.5 et $-0.5$. C'est un changement important, et un avertissement sur les données d'entrée.

Pour une corrélation égale à 1, la matrice est singulière et le facteur a un pivot nul. Il faut alors utiliser une seule normale pour les deux variables, ou tirer à partir des vecteurs propres, avec $\sqrt{\lambda}$.

\iqwrong{«~Ajouter un petit nombre à la diagonale.~» Cela sauve une matrice qui frôle la validité, et ne fait rien pour une valeur propre de $-0.8$ tant que les corrélations ne sont pas diluées au point d'être inutiles.}
\end{solution}

\begin{solution}{iq:iv3:python-and-pandas:34}
Le \Cref{lst:iv3:pp:mc} donne 10.4585 avec une erreur type de 0.0073, sur un million de paires, soit à 1.1 erreur type de la valeur exacte. Traiter les deux membres de la paire comme indépendants donne ici 0.0104, valeur trop élevée. Sur un payoff presque symétrique, comme un straddle, elle est trop faible, ce qui est la direction dangereuse.
\omcode{code/interviews-3/02-python-and-pandas/python/iv_mc_call.py}{1}{21}{Blocs de paires antithétiques, un échantillon par paire.\protect\hfil\protect\break}\label{lst:iv3:pp:mc}

La moyenne de la paire est l'échantillon indépendant. La variance doit donc provenir de ces moyennes, rapportées au nombre de paires, ce que fait le listing. Pour l'option d'achat, les deux branches, haussière et baissière, sont négativement corrélées~: un $z$ élevé rapporte, et un $z$ faible ne rapporte rien. Moyenner la paire retire donc plus de variance que ne le suppose la formule naïve, et l'erreur naïve est trop grande. Un straddle au forward rapporte des deux côtés, si bien que les membres sont positivement corrélés, et la paire porte moins d'information que deux tirages. L'erreur naïve promet alors plus de précision que vous n'en avez, et le test le confirme sur un tel payoff.

La mémoire reste bornée, car seuls les totaux cumulés survivent à chaque bloc, si bien que davantage de trajectoires ne coûtent que du temps. Avec un payoff dépendant du chemin, la même idée s'applique aux pas. Un tableau complet d'un million de trajectoires et de 252 pas occupe environ 2~Go, alors qu'une boucle sur les pas ne conserve que l'état courant.

Un contrôle exact ne demande aucune formule d'option~: à volatilité nulle, le prix doit valoir $S - K e^{-rT}$ lorsque l'option est dans la monnaie, à une erreur de l'ordre de l'arrondi près. Un second contrôle est statistique~: à strike 0, l'option d'achat vaut le titre, si bien que l'estimation doit se trouver à quelques erreurs types du spot, ce qui teste la dérive.

Juger d'après la proximité avec la formule fermée. Une seule graine prouve peu de chose~: il faut vérifier que l'écart reste voisin d'une erreur type sur de nombreuses graines, et que l'erreur décroît en $1 / \sqrt{n}$.
\end{solution}

\begin{solution}{iq:iv3:python-and-pandas:35}
Avec les ordres 9, $-1$ et $L = 8$, la somme écrêtée donne 8, 8, alors que la position réelle est 8, 7. Un seul ordre ne peut jamais le montrer, si bien que deux est le plus petit nombre possible. Un plafonnement modifie chaque position ultérieure, alors que \texttt{cumsum} garde le total non plafonné. Le \Cref{lst:iv3:pp:position} donne la boucle correcte.
\omcode{code/interviews-3/02-python-and-pandas/python/iv_position.py}{1}{20}{Un plafond modifie tout ce qui le suit.}\label{lst:iv3:pp:position}

Le premier ordre est ramené de 9 à 8, de sorte qu'une action n'est jamais achetée. La somme courante la compte encore, si bien qu'après $-1$ elle affiche 8 alors que la position est 7. Il suffit d'un seul plafonnement et d'un seul renversement.

La version à double plafond dépend du chemin, puisque chaque pas a besoin du résultat précédent~: elle ne se réduit donc pas à une opération sur tableau de la manière habituelle. Il faut garder la boucle, qui est claire et testée, et si elle apparaît dans un profilage, la compiler. C'est une bonne réponse~: la vectorisation n'est qu'un moyen, et une version rapide mais fausse ne fait gagner aucun temps.

Avec un seul plafond supérieur, il existe une forme fermée. Soit $c$ la somme cumulée. Le montant total écrêté jusqu'à présent vaut $\max(0, \max_{j \le i} c_j - L)$, si bien que la position vaut $c_i$ moins ce montant, soit \texttt{capped\_above} dans le listing. Le total écrêté ne croît que lorsque $c$ atteint un nouveau maximum au-dessus de $L$. Le test le compare à la boucle sur des ordres aléatoires, et pour $L$ supérieur à toutes les sommes, il se réduit à $c$ lui-même.

\iqwrong{Écrêter chaque ordre à $\pm L$ avant le \texttt{cumsum}. Cela plafonne l'ordre et non la position, si bien que la somme peut encore dépasser largement $L$.}
\iqfollow{Et avec \texttt{pandas}~?}{\texttt{Series.clip} a le même défaut, puisqu'il agit aussi sur la somme finale.}
\end{solution}

\begin{solution}{iq:iv3:python-and-pandas:36}
Le cache les retient. \texttt{lru\_cache} appliqué à une méthode fait de \texttt{self} une partie de la clé, et sans borne il conserve toutes les clés indéfiniment. Après 50 requêtes, \texttt{cache\_info} indique 0 succès, 50 échecs et une taille de 50. Chaque requête porte un nouveau \texttt{self}, de sorte qu'un résultat stocké ne peut plus jamais être retrouvé.

\texttt{gc.collect()} ne libère rien, car il ne s'agit pas de cycles perdus. Ces objets sont atteignables depuis l'objet fonction défini dans la classe, qui vit aussi longtemps que le processus. Dans notre test, une \texttt{weakref} vers un \texttt{Pricer} terminé est encore vivante après coup.

Pour repérer une croissance, comparer deux instantanés \texttt{tracemalloc} pris à un certain intervalle. Ici, la comparaison sur 50 requêtes attribue 4.0~Mo à la seule ligne qui construit \texttt{curve}, et cela désigne ce qui est conservé. Elle indique où un bloc a été alloué, et non qui le détient~: \texttt{gc.get\_referrers} ou un outil de tas répond alors à cette question.

Elle ne voit que la mémoire demandée par les allocateurs de Python, si bien qu'un appel direct à \texttt{malloc} depuis une bibliothèque C reste invisible. De plus, une baisse du total côté Python ne réduit pas forcément la taille résidente du processus, car la mémoire libérée peut rester entre les mains de l'allocateur. Il faut donc comparer le total tracé à la taille résidente~: si seule la seconde croît, il faut chercher plus bas que Python.

La solution consiste à mettre en cache une fonction de valeurs simples au niveau du module, avec une borne, de sorte que \texttt{self} ne figure pas dans la clé. Ou bien donner à chaque instance son propre dictionnaire, qui disparaît avec elle. Dans le test, l'une ou l'autre version laisse l'instance être libérée.

\texttt{maxsize=128} est une borne et non une correction. La mémoire cesse de croître à 128 pricers, chacun toujours retenu pour rien, et le nombre de succès reste à 0.

\iqwrong{«~Python possède un ramasse-miettes, donc cela ne peut pas fuir.~» Il ne récupère que ce qui n'est pas atteignable.}
\end{solution}

\begin{solution}{iq:iv3:python-and-pandas:37}
Les empreintes des chaînes sont randomisées à chaque processus, si bien que l'ordre d'itération d'un ensemble de chaînes change d'une exécution à l'autre. La première ligne dépend de l'ordre, car les quatre symboles sont à égalité et \texttt{max} retient le premier rencontré. La seconde en dépend aussi, car la somme en virgule flottante absorbe ou non le 1.0 selon le moment où il arrive. Les empreintes des entiers ne sont pas randomisées, si bien que des codes ne feraient pas apparaître ce phénomène.

Dans les tests, six valeurs de \texttt{PYTHONHASHSEED} donnent quatre gagnants, et une somme valant tantôt 0.0, tantôt 1.0. Avec une graine fixée, la sortie se répète. Un ensemble d'entiers donne un seul ordre pour toutes les graines.

La correction consiste à ne plus dépendre de l'ordre d'un ensemble là où l'ordre compte~: il faut itérer sur \texttt{sorted(symbols)} et préciser un critère de départage dans la clé, comme \texttt{(score, symbol)}. Un dictionnaire est sûr, puisqu'il conserve l'ordre d'insertion par règle du langage, à condition que ce qui le remplit arrive dans un ordre défini.

Fixer \texttt{PYTHONHASHSEED} rend les exécutions reproductibles, et cela masque le défaut. L'ordre reste arbitraire, et un autre univers, une mise à jour de Python ou un symbole ajouté modifie le résultat sans aucun changement de code. C'est acceptable comme outil de diagnostic~: si la correction fait concorder les exécutions, on a trouvé la classe de bogues.

À vérifier ensuite~: les travaux parallèles recueillis dans l'ordre de fin d'exécution, les listings de répertoires lus sans tri, l'algèbre linéaire multithread, dont l'arrondi peut varier avec le nombre de threads, et une requête de base de données sans \texttt{ORDER BY}. Il faut trouver la première divergence en journalisant une empreinte de l'état chaque jour dans deux exécutions.

\iqwrong{«~Fixer \texttt{np.random.seed}.~» Cela initialise le générateur de NumPy et n'a aucun rapport avec les empreintes des chaînes ni avec l'ordre des ensembles.}
\iqfollow{Un écart au 15\up{e}~chiffre vaut-il la peine d'être recherché~?}{Oui, car de minuscules écarts peuvent inverser un départage ou un seuil, et ensuite on ne peut plus distinguer une modification du code d'un bruit.}
\end{solution}

\begin{solution}{iq:iv3:python-and-pandas:38}
\Cref{lst:iv3:pp:ic} le calcule à l'aide de rangs par groupe et de sommes par groupe, et il concorde avec \texttt{scipy.stats.spearmanr} date par date dans le test. Pour des variables conjointement gaussiennes de corrélation 0.1, on attend $\frac{6}{\pi} \arcsin(0.05) = 0.0955$. La valeur d'un seul jour avec 500 titres a un bruit voisin de $1 / \sqrt{500} = 0.045$.
\omcode{code/interviews-3/02-python-and-pandas/python/iv_ic.py}{1}{11}{Corrélation de rang par date, sans boucle sur les dates.\protect\hfil\protect\break}\label{lst:iv3:pp:ic}

Une corrélation de rang est une corrélation de Pearson appliquée aux rangs~: il faut donc ranger à l'intérieur de chaque date, retrancher la moyenne de chaque date, puis former la somme des produits divisée par la racine du produit des deux sommes des carrés. Le classement intervient après le \texttt{dropna}, afin que les deux colonnes classent les mêmes lignes, pour la même raison que dans toute statistique par paires.

\texttt{df.price.shift(-1)} sur une table triée par symbole donne à la dernière ligne de chaque symbole le premier cours du symbole suivant. Cela produit un rendement de l'ordre de $-40\%$, sans NaN pour vous prévenir, sur une ligne par symbole. Un décalage par groupe donne NaN à cet endroit, et le \texttt{dropna} le supprime.

Une simulation de 500 titres sur 199 jours donne une moyenne de 0.0955 et un écart quotidien de 0.041, proche de cette estimation, pour une statistique de Student de 33, à savoir la moyenne divisée par l'écart, rapporté à $\sqrt{199}$, soit $0.0955 / (0.041 / \sqrt{199})$. Un signal faible un jour donné devient donc net sur un long échantillon. Aux extrêmes, un signal égal au rendement donne exactement 1, son opposé donne $-1$, et un signal constant donne NaN.

Avec une cible à cinq jours, des jours voisins partagent quatre de leurs cinq rendements~: les valeurs quotidiennes se recouvrent, et la statistique de Student naïve est trop élevée, d'environ $\sqrt{5}$ si le signal évolue lentement. Il faut donc n'échantillonner qu'un jour sur cinq, ou utiliser une erreur qui tienne compte de ce recouvrement.

\iqwrong{Moyenner l'IC sur les titres, à partir d'une série temporelle par titre. Cela répond à une autre question~: cela mêle le niveau propre à chaque titre, et demande bien plus de données par titre.}
\end{solution}

\begin{solution}{iq:iv3:python-and-pandas:39}
Le \Cref{lst:iv3:pp:asofnp} remplace les instants par leurs rangs dans l'ensemble des deux tableaux, construit une clé entière par symbole et rang, puis effectue une seule recherche dichotomique par transaction. Il concorde avec \texttt{pd.merge\_asof} et avec une boucle lente sur des données aléatoires. Un multiplicateur assez grand pour les instants fait déborder \texttt{int64}, dont la limite vaut environ $9.2 \times 10^{18}$.
\omcode{code/interviews-3/02-python-and-pandas/python/iv_asof_numpy.py}{1}{15}{Une jointure as-of par un seul tri et une seule recherche.\protect\hfil\protect\break}\label{lst:iv3:pp:asofnp}

L'idée consiste à réduire deux clés de tri à une seule. Si la clé ordonne d'abord par symbole puis par instant, une recherche dichotomique de la clé d'une transaction trouve la dernière cotation inférieure ou égale à celle-ci, et il ne reste qu'une vérification~: que cette cotation porte le même symbole. Si elle porte un autre symbole, alors ce symbole n'a aucune cotation antérieure ou égale à la transaction, car une telle cotation aurait une clé comprise entre les deux.

Avec $10^{19}$ lui-même, NumPy 2.2 lève \texttt{OverflowError}, car cette valeur ne tient pas dans \texttt{int64}. Un multiplicateur qui tient doit tout de même dépasser $1.8 \times 10^{18}$ pour les instants, si bien que le symbole 5 suivi d'un instant, soit environ $1.08 \times 10^{19}$, dépasse déjà la limite, et NumPy reboucle sans erreur, ce qui fausse l'ordre. \texttt{np.unique} avec \texttt{return\_inverse} règle le problème. Il y a au plus $m + n$ instants distincts, de sorte que la clé reste inférieure à $5000 (m + n)$, et les rangs conservent l'ordre, ce qui suffit à la recherche. L'énoncé ne donne aucun type pour les codes~: le listing les convertit donc d'abord en \texttt{int64}, car des codes en \texttt{int32} reboucleraient sans erreur dès que $m + n$ atteint environ un million, et des codes en \texttt{int16} lèvent une erreur.

\texttt{side=\textquotesingle right\textquotesingle} trouve les cotations antérieures ou égales à la transaction, ce qu'exige l'énoncé. Avec \texttt{left}, il ne retiendrait que les cotations strictement antérieures. Le tri stable argsort conserve l'ordre d'arrivée entre clés égales, de sorte que la recherche tombe sur la dernière arrivée.

Ce sont les tris qui dominent, en $O((m + n) \log (m + n))$, et les recherches coûtent $O(m \log n)$. Si les deux entrées arrivent déjà triées, un parcours de fusion par symbole est linéaire et demande en plus moins de mémoire. Ici, l'étape de rangs nécessite un tableau de $m + n$ instants, et un tableau de rangs.

\iqwrong{\texttt{np.searchsorted} sur \texttt{quote\_t} seul, puis un parcours en arrière jusqu'au bon symbole. Sur une table chargée, ce parcours peut traverser de nombreuses cotations d'autres symboles, si bien que son coût par transaction n'est pas borné.}
\end{solution}

\begin{solution}{iq:iv3:python-and-pandas:40}
La méthode des sommes cumulées donne un grand nombre négatif à la ligne 999, $-77$ sur notre tirage, alors que la pente vraie vaut 2.00. Le \Cref{lst:iv3:pp:rollbeta} l'évite en $O(n)$, car la somme de chaque fenêtre n'utilise que les valeurs situées à l'intérieur de cette fenêtre. La variance glissante de pandas 2.3.3 affiche là un multiple entier de $128/49 = 2.61$, soit $-2.6$ sur notre tirage et positive ou nulle sur d'autres, alors que la valeur vraie vaut environ $10^{-4}$. Son \texttt{rolling.corr} affiche NaN, un nombre minuscule, ou une valeur supérieure à 1.
\omcode{code/interviews-3/02-python-and-pandas/python/iv_rolling_beta.py}{1}{26}{Sommes glissantes qui oublient ce qui a quitté la fenêtre.\protect\hfil\protect\break}\label{lst:iv3:pp:rollbeta}

Masquer d'abord les deux colonnes fait que les cinq sommes glissantes portent sur les mêmes lignes, et la pente vaut $(n S_{xy} - S_x S_y) / (n S_{xx} - S_x^2)$.

Le dégât vient de $x^2 = 10^{18}$. Les doubles de cette taille sont espacés de 128, si bien que, dès que ce terme entre dans la somme courante, chaque petit terme suivant est arrondi à zéro. La soustraction de deux telles sommes donne exactement 0, alors que la vraie valeur de $S_{xx}$ est voisine de $5 \times 10^{-3}$. Le dénominateur vaut alors $-S_x^2$, ce qui rend la pente grande et négative~: de $-16$ à $-44\,000$ sur dix tirages. La faute persiste dans chaque ligne ultérieure, longtemps après que la mauvaise valeur a quitté la fenêtre.

Le listing découpe la série en blocs de longueur $w$. Toute fenêtre de longueur $w$ est alors un suffixe d'un bloc suivi d'un préfixe du bloc suivant, et ces deux parties sont de simples sommes cumulées à l'intérieur d'un bloc. Une somme ne contient donc jamais une valeur extérieure à sa fenêtre, et chaque ligne coûte une seule addition. Dans le test, il concorde avec un ajustement lent sur chaque fenêtre à l'écart de la mauvaise valeur. Les fenêtres qui contiennent cette valeur restent dominées par elle, comme l'exige l'arithmétique, si bien que le nettoyage reste une étape distincte.

Les chiffres de pandas sont ce que fait la version 2.3.3 ici. Sa variance glissante ajoute et retire les valeurs au fil du déplacement de la fenêtre, et elle laisse un résidu de cette étape sur cette entrée, alors que \texttt{rolling.cov} sur ces mêmes colonnes reste juste. Il faut effectuer un contrôle de ce type sur sa propre version de la bibliothèque, car le comportement peut changer.

\iqwrong{«~Utiliser une somme compensée pour corriger la dérive.~» Cela réduit l'erreur, mais la mauvaise valeur passe tout de même par un état courant. Repartir à zéro à chaque bloc supprime cette dépendance.}
\end{solution}
